\scanpage{24}
d'où
\[
\unsure{\partial_n|\delta_n}:\underline H^1_Y(G_{mX})\longrightarrow\underline H^2_Y(\mu_{nX})
\]
et nous posons
\[
\deleted{\unsure{\varphi}}\ \underline{\unsure{\gamma|\sigma}}_{Y/X}^{(n)}=\unsure{\partial_n|\delta_n}(\underline\beta_{Y/X})\in\Gamma(X,\underline H^2_Y(\mu_{nX}))
\]
en laissant tomber les exposants $n$ s'il n'y a pas d'ambiguïté possible.

\inserted{La définition de} Cette section \deleted{satisfait} commute à la localité de Zariski (et $=$ à tt changement de base $X'\to X$ tel que $Y'$ reste une ``int. complète'' i.e. $\varphi'$ reste section régulière) ce qui permet de définir $\gamma_{Y/X}^{(n)}$ en général (sans équation globale). Si \inserted{maintenant} $Y$ est intersection complète de codim $d$, alors \unsure{localt} $Y$ est de la forme
\[
Y=Y_1\cap\ldots\cap Y_d
\]
et on posera
\[
\underline{\unsure{\gamma|\sigma}}_{Y/X}
=\underline{\unsure{\gamma|\sigma}}_{Y_1/X}^{\deleted{\unsure{(n)}}}\ .\ \ldots\ .\ \underline{\unsure{\gamma|\sigma}}_{Y_d/X}
\in\Gamma(Y,\underline H^{2d}_Y(\mu_n^{\otimes d}))
\]
(\uline{cup-produit}), et on prouve que cette section ne dépend pas du choix des $Y_i$, donc se définit globalement.

Propriétés générales
\mspage{23}
